% =====================================================================
% LAMPIRAN E: KONFIGURASI RUNPOD DAN WORKFLOW SSH
% =====================================================================
\chapter{Konfigurasi RunPod, Workflow SSH, dan Manajemen Lingkungan}
\label{app:runpod-workflow}

Lampiran ini merinci alur kerja akses \textit{pod} RunPod via SSH serta praktik manajemen lingkungan yang diringkas pada Bagian~\ref{sec:runpod}. Konteks visual alur kerja sudah disajikan pada Gambar~\ref{fig:kerangka-pelatihan-setup} di Bab~\ref{ch:metode}.

\section*{E.1 Akses Pod via SSH dan \textit{Workflow}}

Akses ke \textit{pod} dilakukan melalui kunci publik SSH yang didaftarkan pada akun RunPod. \textit{Workflow} pelatihan terdiri atas langkah-langkah berikut.

\begin{enumerate}
    \item \textbf{Penyiapan kunci SSH}: pembuatan pasangan kunci pada mesin lokal dan unggah kunci publik ke profil RunPod.
    \item \textbf{Login ke pod}: \texttt{ssh root@<pod-host>\allowbreak{} -p <port>\allowbreak{} -i \textasciitilde/.ssh/\allowbreak{}id\_ed25519}.
    \item \textbf{Sinkronisasi dataset}: pengiriman dataset bersih dan direktori gambar dari mesin lokal ke \textit{pod} menggunakan \texttt{rsync} melalui kanal SSH.
    \item \textbf{Penyiapan lingkungan Python}: pemasangan dependensi pelatihan, mencakup \texttt{torch}, \texttt{transformers}, \texttt{timm}, \texttt{catboost}, dan \texttt{onnxruntime-gpu}.
    \item \textbf{Eksekusi pelatihan}: peluncuran Jupyter Lab pada \textit{pod} dan akses dari peramban lokal melalui \textit{SSH tunnel} ke \textit{port} 8888.
    \item \textbf{Penyimpanan \textit{artifact}}: \textit{embedding} dan model akhir disimpan pada \textit{persistent volume} \textit{pod}.
    \item \textbf{Pengunduhan hasil}: \textit{embedding cache} dan \textit{checkpoint} CatBoost diunduh kembali ke mesin lokal menggunakan \texttt{rsync} setelah pelatihan selesai.
\end{enumerate}

\section*{E.2 Manajemen Lingkungan}

Daftar dependensi dengan versi presisi dipertahankan sebagai berkas konfigurasi pelatihan sehingga lingkungan dapat direproduksi pada \textit{pod} baru. Eksekusi \textit{forward pass} \textit{encoder} memanfaatkan presisi campuran (\textit{Automatic Mixed Precision}) FP16 untuk mengurangi konsumsi VRAM dan mempercepat komputasi. Ukuran \textit{batch} disesuaikan per \textit{encoder}, contohnya 32 untuk DINOv3 pada masukan 224x224 dan 8 untuk EVA-02 pada masukan 448x448, berdasarkan ketersediaan VRAM.
